\chapter{Cadre général du projet}
\label{chap:cadre-general}
\begin{spacing}{1.1}
\minitoc
\thispagestyle{MyStyle}
\end{spacing}
\newpage

\section{Introduction}
Ce chapitre introduit le cadre institutionnel, organisationnel et technique du projet. Il précise le besoin métier, les objectifs poursuivis, les acteurs impliqués ainsi que le périmètre de réalisation retenu.

\section{Présentation de l'organisme d'accueil}
\subsection{DECADE Tunisie et contexte d'activité}
DECADE Tunisie est une filiale spécialisée dans l'ingénierie e-commerce. Le stage s'est déroulé au sein de l'entité de Sfax, dans un contexte de projets orientés qualité de données, performance de recherche et optimisation des parcours utilisateurs.

L'environnement technique de référence mobilise notamment \textbf{SAP Commerce Cloud (Hybris)} pour la gestion du catalogue, \textbf{Apache Solr} pour la recherche et les facettes, ainsi qu'une architecture orientée \textbf{API REST} pour l'interopérabilité des services.

\subsection{Cadre académique du stage}
Ce travail s'inscrit dans le cadre du projet de fin d'études de la filière Génie Informatique et Mathématiques Appliquées de l'ENIS, sur la période de février à juillet 2026. Il vise à produire une contribution concrète, exploitable et argumentée scientifiquement.

\section{Problématique}
Dans un catalogue e-commerce volumineux, la qualité des fiches produit conditionne directement la pertinence de la recherche. Or, l'enrichissement manuel des attributs présente plusieurs limites:

\begin{itemize}
  \item forte charge opérationnelle pour les équipes métier;
  \item hétérogénéité de saisie (valeurs redondantes, variantes lexicales, omissions);
  \item retard d'indexation de données utiles pour Solr;
  \item impact négatif sur les filtres dynamiques et le ranking.
\end{itemize}

La problématique centrale peut être formulée ainsi:\\
\textit{Comment automatiser l'enrichissement des attributs produit à partir de données multimodales, tout en assurant la qualité, la traçabilité et l'intégration native dans l'écosystème Hybris/Solr ?}

\section{Objectifs du projet}
\subsection{Objectif principal}
Concevoir une solution d'enrichissement automatique des données produit par IA multimodale, intégrée à SAP Commerce Cloud et exploitable dans Apache Solr pour améliorer la recherche e-commerce.

\subsection{Objectifs spécifiques}
\begin{itemize}
  \item Extraire automatiquement des attributs structurés (couleur, matière, catégorie, motif, saisonnalité, description);
  \item Mettre en place un espace de validation humaine (validation, correction, rejet) pour maîtriser la qualité;
  \item Injecter les attributs validés dans Hybris avec une logique de mapping et de gestion de statuts;
  \item Alimenter Solr pour améliorer facettes, filtres et pertinence des résultats;
  \item Exposer les données enrichies via des services API pour faciliter la réutilisation.
\end{itemize}

\subsection{Indicateurs cibles}
Afin d'évaluer l'apport de la solution, les indicateurs suivants sont retenus:
\begin{itemize}
  \item taux de complétion des attributs avant/après enrichissement;
  \item part des suggestions validées manuellement;
  \item score de confiance moyen par attribut;
  \item impact sur la qualité de recherche (facettes disponibles, précision perçue).
\end{itemize}

\section{Acteurs et interactions principales}
Le système met en jeu des acteurs humains et techniques complémentaires:
\begin{itemize}
  \item \textbf{Gestionnaire produit}: valide, corrige ou rejette les attributs générés;
  \item \textbf{Administrateur Hybris}: supervise le mapping et l'import des données validées;
  \item \textbf{Service IA multimodal}: propose les enrichissements avec score de confiance;
  \item \textbf{Apache Solr}: indexe les nouveaux attributs pour améliorer la recherche;
  \item \textbf{Client final}: bénéficie indirectement d'une navigation plus pertinente.
\end{itemize}

\section{Périmètre et contraintes}
\subsection{Périmètre fonctionnel}
Le périmètre du PFE couvre les briques suivantes:
\begin{itemize}
  \item enrichissement automatique d'attributs prioritaires;
  \item validation humaine via interface dédiée;
  \item intégration Hybris des attributs validés;
  \item indexation Solr des nouveaux champs;
  \item exposition des données enrichies via API.
\end{itemize}

\subsection{Contraintes de réalisation}
Les contraintes majeures identifiées sont:
\begin{itemize}
  \item compatibilité avec le modèle de données Hybris existant;
  \item maîtrise de la latence et du coût d'inférence LLM en traitement batch;
  \item robustesse des échanges API (gestion d'erreurs, reprise, traçabilité);
  \item qualité éditoriale multilingue et cohérence des attributs générés.
\end{itemize}

\section{Démarche méthodologique et planification}
Le projet suit une démarche progressive structurée en phases:
\begin{itemize}
  \item cadrage et étude de l'existant;
  \item spécifications fonctionnelles et techniques;
  \item conception d'architecture;
  \item implémentation incrémentale par sprints;
  \item tests, mesure des indicateurs et consolidation documentaire.
\end{itemize}

Cette organisation permet d'aligner la livraison technique avec les objectifs métier, tout en conservant une capacité d'ajustement continue.

\section{Conclusion}
Ce chapitre a posé les fondations du projet: contexte entreprise, problématique, objectifs, acteurs, périmètre et contraintes. Il constitue le référentiel de cadrage sur lequel s'appuient les choix d'analyse et de conception présentés dans la suite du rapport.
